\section*{Chapter \ref{ch:b3:project-replicated-store} --- Project: A Replicated Key-Value Store}

\begin{solution}{exo:b3:project-replicated-store:1}
The client sends \code{SET ferry 8h15} followed by a newline: $3 + 1 + 5 + 1 + 4 + 1 = 15$ bytes. The replies: \code{SET a=b 1} gets \code{ERR a key may not contain =}; \code{SET a b=c} gets \code{OK}, since only keys are restricted (\code{DUMP} prints \code{a=b=c}, and a reader that splits at the first \code{=} gets the pair back); \code{GET a b} gets \code{ERR GET takes one key}; \code{DEL} with two spaces before \code{a} is \code{DEL a}, because the parser drops empty words, and gets \code{OK} here, since \code{a} was just set; \code{status} gets \code{ERR unknown command status}, because commands are case-sensitive. The three bytes \code{GET} wait in the server's buffer: \code{read_line} cannot know whether more of the line will come until a newline or the end of the stream arrives. When the client closes, \code{read_line} returns the partial line, which is parsed and answered with \code{ERR GET takes one key}, written to a connection nobody reads. The test \code{exercise_1_lines} checks each case.
\end{solution}

\begin{solution}{exo:b3:project-replicated-store:2}
The records are 21, 29, 40 and 42 bytes long and end at offsets 21, 50, 90 and 132. At 100 bytes, three records are whole: the vote, the \code{NOOP} and \code{SET bridge open}; the fourth has its 8-byte header and 2 bytes of its 34-byte payload. By \cref{prop:b3:project-replicated-store:replay} the replay returns the three, and \code{Wal::open} cuts the file to 90 bytes. A client reads \code{VALUE open}. The client of \code{SET bridge open} may have received \code{OK}: its record was forced before the next one was written, and the reply was sent just after; whether it left before the crash depends on the timing. The client of \code{SET bridge closed} cannot have received \code{OK}, which is sent only after \code{sync_all} returns on the whole record. The test \code{exercise_2_a_crash_after_100_bytes} checks the sizes and the replay.
\end{solution}

\begin{solution}{exo:b3:project-replicated-store:3}
A voter compares the pair (\omterm{def:b3:replication-and-consensus:raft}{election term} of the last entry, length): S1 has $(2, 3)$, S2 has $(3, 3)$, S3 has $(1, 2)$. A candidate needs two votes, its own and one more. S1 gets S3's vote, since $(2, 3) \geq (1, 2)$, and not S2's: it can win. S2 gets both: it can win. S3 gets neither: it cannot. If S1 wins \omterm{def:b3:replication-and-consensus:raft}{election term}~4, its first append to S2 carries its third entry, of \omterm{def:b3:replication-and-consensus:raft}{election term}~2; S2 finds \omterm{def:b3:replication-and-consensus:raft}{election term}~3 at that index, removes its entry and takes S1's. That entry of \omterm{def:b3:replication-and-consensus:raft}{election term}~3 was held by S2 alone, one replica out of three, so it was never committed, and no client was told \code{OK} for it. The test \code{exercise_3_who_can_lead} replays the six requests for a vote and finds the winners S1 and S2.
\end{solution}

\begin{solution}{exo:b3:project-replicated-store:4}
\omoutput{code/bachelor-3/32-project-replicated-store/outputs/ex4_session.txt}{The session of the exercise, captured.}\label{out:b3:project-replicated-store:ex4}
With replica~2 dead, replicas 1 and 3 are still a majority, so \code{SET y 2} is redirected and committed. With replica~3 dead too, replica~1 is alone. It never hears a larger \omterm{def:b3:replication-and-consensus:raft}{election term}, so it still believes it leads: it appends \code{SET z 3} and forces it, but no other replica answers, and the entry is never committed. The client gives up after its read timeout of 2 seconds and prints \code{no answer}; the write stays pending, and the client cannot tell whether it will ever take effect. The read is still answered from the map, and \code{STATUS} still says \code{leader}: a replica cannot tell that it is alone from what it receives, only from what it does not receive.
\end{solution}

\begin{solution}{exo:b3:project-replicated-store:5}
The file holds two records of $8 + 17 + 15 = 40$ and $8 + 17 + 17 = 42$ bytes, 82 in all; cut at 77, the second record has its header, whose length says 34, but only 29 bytes of payload. The trusting version takes those 29 bytes as the payload, so the command is its last $29 - 17 = 12$ bytes, \code{SET bridge c}: the replica replays a write that no client ever made, and \code{GET bridge} returns \code{c}. The fix is the code of \cref{lst:b3:project-replicated-store:decode}: take the payload only if all \code{len} bytes are there, and accept it only if its checksum matches. The test \code{ex5_torn_test.rs} shows the invented value with the trusting version, then cuts the file at each of its 83 positions and checks that the fixed decoder returns a prefix of the records and the length of their encodings; a flipped bit in the last command, which no length check can see, stops the replay at that record.
\end{solution}

\begin{solution}{exo:b3:project-replicated-store:6}
When it applies it. On receipt, the leader's map does not yet contain the writes before the \code{CAS} in the log that are not committed yet, so the comparison would be made against the wrong state; and the followers, which see only the log, must reach the same answer from it. So \code{CAS} is proposed like \code{SET}, and its result is computed by \code{apply}, from the map as it is just before the entry, in log order:
\begin{itemize}
\item value equal to \code{old}: set \code{new}, reply \code{OK};
\item another value $v$: reply \code{VALUE v};
\item no value: reply \code{NOTFOUND}.
\end{itemize}
The replica that proposed it sends that reply to the client, as for any write. The test \code{ex6_cas_test.rs} applies the log \code{SET seat free}, \code{CAS seat free ada}, \code{CAS seat free bob}, \code{CAS desk free ada} on three replicas, which all answer \code{OK}, \code{OK}, \code{VALUE ada}, \code{NOTFOUND} and end with \code{seat} holding \code{ada}; a \code{CAS} with two words is refused by the parser.
\end{solution}

\begin{solution}{exo:b3:project-replicated-store:7}
\emph{A leader cut off.} Replica~1 leads \omterm{def:b3:replication-and-consensus:raft}{election term}~1, and \code{a} is 1. The network cuts replica~1 off; replicas 2 and 3 elect replica~2 in \omterm{def:b3:replication-and-consensus:raft}{election term}~2, which commits \code{SET a 2} and replies \code{OK}. A client then sends \code{GET a} to replica~1, which has heard of no larger \omterm{def:b3:replication-and-consensus:raft}{election term}, still leads in its own eyes, and answers \code{VALUE 1}. The read started after a write of 2 had returned, so every linearisation places it after that write, and it should return 2: the history is not \omterm{def:b3:concurrency:linearisability}{linearisable}. The test \code{a_deposed_leader_reads_stale} builds this run.

\emph{A leader just elected.} Replica~1 commits \code{SET a 2} with replica~2, replies \code{OK}, and crashes before its next append tells replica~2 that the entry is committed. Replica~2 is elected (its log is as up to date as replica~3's or more), but its \omterm{def:b3:replication-and-consensus:raft}{commit index} is still below that entry, so its map still has the old value, and a \code{GET a} answered at once returns it.

\emph{A read index.} On a \code{GET}, the leader waits until an entry of its own \omterm{def:b3:replication-and-consensus:raft}{election term} is committed, its \code{NOOP}, then records its \omterm{def:b3:replication-and-consensus:raft}{commit index} $c$ as the read's index. It sends an append to every follower and waits for a majority to answer it in its current \omterm{def:b3:replication-and-consensus:raft}{election term}, which shows that no other leader had been elected when the read arrived. It then waits until it has applied index $c$, and answers from its map. Every write acknowledged before the read arrived is committed, so its index is at most $c$, and the answer includes it. In the first run, replica~1's appends reach nobody, and the \code{GET} gets no answer; the client retries elsewhere and reaches replica~2. In the second, replica~2 answers only after its \code{NOOP} is committed, which commits \code{SET a 2} with it, and it returns 2. This is the check of the \omterm{def:b3:replication-and-consensus:raft}{Raft} paper.
\end{solution}

\begin{solution}{exo:b3:project-replicated-store:8}
The leader forces the entry and sends one append to each follower: 2 messages. Each follower forces it and answers: 2 more. The first answer already gives a majority, the leader and one follower, so the leader commits and applies the entry: 4 messages, of which the leader needed 3 before committing. The followers learn that the entry is committed from the \omterm{def:b3:replication-and-consensus:raft}{commit index} carried by the leader's next append, a heartbeat or the next write, so this costs no message of its own. The test \code{messages_per_write} counts the messages, and finds the leader's \omterm{def:b3:replication-and-consensus:raft}{commit index} at 2 and the followers' still at 1. With $N$ replicas: $N - 1$ appends and $N - 1$ answers, $2(N - 1)$ messages, and the leader commits after $\lfloor N/2 \rfloor$ answers. Add the client's request and reply: 2 more.
\end{solution}

\begin{solution}{exo:b3:project-replicated-store:9}
Replay the records into the state they describe, the \omterm{def:b3:replication-and-consensus:raft}{election term}, the vote and the log, then write that state back: one vote record, then one entry record per entry of the log, with indices $1, 2, \ldots$. Replaying the result gives the same \omterm{def:b3:replication-and-consensus:raft}{election term} and vote (the only vote record) and the same log (each entry record extends the log by one). In the test's log, the empty entry holds index 1, the last term writes indices 40 to 44, and the entries at indices 2 to 39 are the last ones written there, so the log has 44 entries and the compacted file 45 records out of 121, less than half the bytes. The test \code{ex9_compact_test.rs} checks that both replay to the same state. In a replica the compacted file must replace the old one in a single step, so that a crash at any moment leaves one of the two files whole, never a mix.
\end{solution}

\begin{solution}{exo:b3:project-replicated-store:10}
Facts used: \cref{def:b3:concurrency:linearisability}; an operation without a reply has no end, so no operation follows it in real time. Call a \emph{solution} of a set $R$ of remaining operations from a value $v$ an order of a subset of $R$ that contains every operation of $R$ with a reply, keeps the real-time order, and is a correct register history starting from $v$. The search is called with all the operations and the empty register, and must return true exactly when a solution exists. (Reads without a reply were removed before: they change nothing, and their result is unknown.)

\emph{The first operation.} If $o$ comes first in a solution, then $\mathrm{start}(o) \leq \mathrm{end}(p)$ for every $p \in R$ with a reply: otherwise $p$ ended before $o$ started, must come before it, and is in the solution, since it has a reply. Conversely, if $o$ satisfies this, placing $o$ first breaks no real-time constraint with any operation of $R$. So the operations the search tries, those not skipped, are exactly the possible first ones.

\emph{Without the memory}, by induction on $|R|$. If every operation of $R$ lacks a reply, the empty order is a solution and the search returns true. Otherwise every solution is non-empty. If the search returns true, it chose a possible first $o$ whose \code{step} from $v$ succeeded, giving $v'$, and found a solution of $R \setminus \{o\}$ from $v'$; putting $o$ in front gives a solution of $R$ from $v$. If a solution exists, let $o$ be its first operation: it is a possible first, \code{step} succeeds on it, and the rest is a solution of $R \setminus \{o\}$ from $v'$, so by induction the call on $o$ returns true, and the search does too.

\emph{The memory.} The result of a call depends only on the pair (remaining set, value), since the operations are fixed. Each call removes one operation, so a pair cannot reappear inside its own call. When a pair is met a second time, its first call has returned: if it had returned true, the whole search would have returned true at once, and the second meeting would not occur; so it returned false, and returning false again is correct.
\end{solution}

\begin{solution}{exo:b3:project-replicated-store:11}
\emph{Lost vote records.} The vote records hold the \omterm{def:b3:replication-and-consensus:raft}{election term} as well as the vote, so the restarted node is back in \omterm{def:b3:replication-and-consensus:raft}{election term}~0 with no vote. It can vote a second time in an \omterm{def:b3:replication-and-consensus:raft}{election term} in which it has already voted, and two candidates can then each count a majority in the same \omterm{def:b3:replication-and-consensus:raft}{election term}. A short run: node~1 leads \omterm{def:b3:replication-and-consensus:raft}{election term}~1 and writes entries; all three nodes crash and restart without their vote records; whichever wins the next election leads \omterm{def:b3:replication-and-consensus:raft}{election term}~1 again, with a log that may differ from the first leadership's at the same indices, so the basis of log matching, one entry per index and \omterm{def:b3:replication-and-consensus:raft}{election term}, is gone. Election safety is the property that the checker sees break, and it breaks first in 106 runs out of 107. Every election after such restarts reuses \omterm{def:b3:replication-and-consensus:raft}{election term} numbers that were already used, and the checker notices at the first leader elected; the other properties need a later, more particular sequence of events.

\emph{Records lost after sending.} A follower appends an entry, sends its success, and crashes before the record is on its disk. The leader counts it, commits the entry and applies it, then crashes. The follower restarts without the entry, and it and the third node, which never received it, form a majority: one of them is elected with the other's vote, and the new leader lacks an applied entry. That is leader completeness, which breaks first in 17 runs out of 18, and then state machine safety, when the new leader's own entry is applied at the same index. Followers write an entry record in almost every step, and a vote record only at elections, so the lost records are almost always entries. The store's main loop rules this fault out by forcing \code{out.persist} before the first message leaves (\cref{lst:b3:project-replicated-store:apply}). The test \code{faults_are_caught} checks the counts and the first seed of each fault, 21 and 25, and that the correct code breaks nothing in the same 2,000 runs.
\end{solution}

\begin{solution}{exo:b3:project-replicated-store:12}
With one replica and no batching, each write waits for its own flush, and the main loop does one at a time. The time of a flush is then about $1/257$ s, so $f \approx 3.9$ ms; the rest of the work, parsing and loopback messages, is small beside it.

\emph{The model.} If the loop finds $b$ writes waiting each time it takes its events, it acknowledges $b$ writes per flush, about $b/f$ per second. With batching, the 1,661 writes per second give $b \approx 1661 \times 0.0039 \approx 6.5$. The gain is bounded by the number of clients: each client has at most one write outstanding, so $b \leq 16$, and the throughput is at most $16/f \approx 4{,}100$ writes per second. It stays far below, because a client is not waiting in the queue while its reply travels back and its next request travels in. The same count follows from Little's law: with 16 clients and 1,661 writes per second, a write takes $16/1661 \approx 9.6$ ms, about $2.5f$, from the client's sending to its receiving.

\emph{Three replicas.} Each write needs two flushes in series, the leader's and then a follower's, plus two message hops. Without batching, the model predicts at least $1/(2f) \approx 128$ writes per second. The measured 154 is above it, because the leader forces the next entry while a follower forces the previous one. With batching, a write takes $16/408 \approx 39$ ms, about four times the time with one replica, where the two flushes in series only explain twice. The rest comes from the setup: the three processes share one disk, so the followers' flushes and the leader's wait for one another. On three machines each flush would have a disk of its own. The bench does not record the flush times of each process, so this explanation is not measured here.
\end{solution}
